% GENERATED FILE. Edit canonical lessons, then run npm run generate:books.
% canonical-source-hash: fnv1a64:489fc664110416f3
% canonical-lessons: GE-C445-schuster, GE-C445-besohlen, GE-C445-polieren, GE-C445-schluesseldienst, GE-C445-aufschliessen, GE-R445-review-438-445

\chapter{Schuhe und Schlüssel}
\label{ch:ge-schuhe-und-schlussel}
\hlchaptermodality{445}

\begin{chapteropening}
\textbf{By the end of this chapter:} \emph{I can get shoes and locks seen to: the cobbler, resoling, polishing, the locksmith and unlocking.}

Complete the last lesson of chapter 445: I can get shoes and locks seen to: the cobbler, resoling, polishing, the locksmith and unlocking.
\end{chapteropening}

\section[der Schuster]{der Schuster — the cobbler, the shoemaker}

\label{lesson:GE-C445-schuster}



\begin{quote}
Before the new one: say the German for the flower shop, then the German for to stick.
\end{quote}

\subsection*{You'll want to know: der Schuster}
\textbf{der Schuster} — \textquotedblleft{}the cobbler, the shoemaker\textquotedblright{}.

\begin{quote}
Der Schuster repariert meine alten Stiefel.
\end{quote}

\begin{quote}
\emph{The cobbler repairs my old boots.}
\end{quote}

Same form in the plural. Also \emph{der Schuhmacher}.

\begin{cousinweb}
\emph{Schuster} comes from Middle High German \emph{schuochsūtære} 'shoe-stitcher' (EtymDB), built from \emph{Schuh} 'shoe' and \emph{sūtære}, from Latin \emph{sūtor} 'cobbler', from \emph{suere} 'to sew' (Lewis \& Short).
\end{cousinweb}

\begin{grammarlens}[title={What you've built}]
One more word for telling, explaining and describing at B1.
\end{grammarlens}

\subsection*{Your turn}
\begin{itemize}
\raggedright
  \item \emph{Say it:} \emph{der Schuster}
  \item \emph{Say it:} \emph{der Schuster}, once more
  \item \emph{Say it:} a sentence of your own with \emph{der Schuster}
  \item \emph{Recall:} say the German for the flower shop, then the German for to stick, then say \emph{der Schuster} again
\end{itemize}

\begin{tcolorbox}[breakable,colback=teal!4,colframe=teal!35!black,title={Before you move on}]
What does der Schuster mean? (\textquotedblleft{}The cobbler, the shoemaker\textquotedblright{}.) Say it once more.
\end{tcolorbox}
\section[besohlen]{besohlen — to resole}

\label{lesson:GE-C445-besohlen}



\begin{quote}
Before the new one: say the German for to stick, then the German for the cobbler.
\end{quote}

\subsection*{You'll want to know: besohlen}
\textbf{besohlen} — \textquotedblleft{}to resole\textquotedblright{}.

\begin{quote}
Der Schuster besohlt die Stiefel bis Freitag.
\end{quote}

\begin{quote}
\emph{The cobbler will resole the boots by Friday.}
\end{quote}

Inseparable: \emph{er besohlt}, \emph{er hat besohlt}.

\begin{cousinweb}
\emph{besohlen} is built from \emph{be-} and \emph{Sohle} 'sole', which goes back to Old High German \emph{sola} (EtymDB) and is usually traced to Latin \emph{solea} 'sandal' (Lewis \& Short).
\end{cousinweb}

\begin{grammarlens}[title={What you've built}]
One more word for telling, explaining and describing at B1.
\end{grammarlens}

\subsection*{Your turn}
\begin{itemize}
\raggedright
  \item \emph{Say it:} \emph{besohlen}
  \item \emph{Say it:} \emph{besohlen}, once more
  \item \emph{Say it:} a sentence of your own with \emph{besohlen}
  \item \emph{Recall:} say the German for to stick, then the German for the cobbler, then say \emph{besohlen} again
\end{itemize}

\begin{tcolorbox}[breakable,colback=teal!4,colframe=teal!35!black,title={Before you move on}]
What does besohlen mean? (\textquotedblleft{}To resole\textquotedblright{}.) Say it once more.
\end{tcolorbox}
\section[polieren]{polieren — to polish}

\label{lesson:GE-C445-polieren}



\begin{quote}
Before the new one: say the German for the cobbler, then the German for to resole.
\end{quote}

\subsection*{You'll want to know: polieren}
\textbf{polieren} — \textquotedblleft{}to polish\textquotedblright{}.

\begin{quote}
Vor der Hochzeit poliert er seine Schuhe.
\end{quote}

\begin{quote}
\emph{Before the wedding he polishes his shoes.}
\end{quote}

Past participle without \emph{ge-}: \emph{ich habe poliert}.

\begin{cousinweb}
\emph{polieren} is inherited from Middle High German \emph{polieren} (EtymDB), a borrowing through Old French \emph{polir} from Latin \emph{polīre} 'to make smooth, to polish' (Lewis \& Short).
\end{cousinweb}

\begin{grammarlens}[title={What you've built}]
One more word for telling, explaining and describing at B1.
\end{grammarlens}

\subsection*{Your turn}
\begin{itemize}
\raggedright
  \item \emph{Say it:} \emph{polieren}
  \item \emph{Say it:} \emph{polieren}, once more
  \item \emph{Say it:} a sentence of your own with \emph{polieren}
  \item \emph{Recall:} say the German for the cobbler, then the German for to resole, then say \emph{polieren} again
\end{itemize}

\begin{tcolorbox}[breakable,colback=teal!4,colframe=teal!35!black,title={Before you move on}]
What does polieren mean? (\textquotedblleft{}To polish\textquotedblright{}.) Say it once more.
\end{tcolorbox}
\section[der Schlüsseldienst]{der Schlüsseldienst — the locksmith (service)}

\label{lesson:GE-C445-schluesseldienst}



\begin{quote}
Before the new one: say the German for to resole, then the German for to polish.
\end{quote}

\subsection*{You'll want to know: der Schlüsseldienst}
\textbf{der Schlüsseldienst} — \textquotedblleft{}the locksmith (service)\textquotedblright{}.

\begin{quote}
Sie hat ihren Schlüssel verloren und ruft den Schlüsseldienst.
\end{quote}

\begin{quote}
\emph{She has lost her key and calls the locksmith.}
\end{quote}

Plural: \emph{die Schlüsseldienste}. The craftsman is \emph{der Schlosser}.

\begin{cousinweb}
\emph{Schlüsseldienst} is built from \emph{Schlüssel} 'key', inherited from Old High German \emph{sluzzil}, from \emph{schließen} 'to lock' (EtymDB), and \emph{Dienst} 'service', from Middle High German \emph{dienst} (EtymDB).
\end{cousinweb}

\begin{grammarlens}[title={What you've built}]
One more word for telling, explaining and describing at B1.
\end{grammarlens}

\subsection*{Your turn}
\begin{itemize}
\raggedright
  \item \emph{Say it:} \emph{der Schlüsseldienst}
  \item \emph{Say it:} \emph{der Schlüsseldienst}, once more
  \item \emph{Say it:} a sentence of your own with \emph{der Schlüsseldienst}
  \item \emph{Recall:} say the German for to resole, then the German for to polish, then say \emph{der Schlüsseldienst} again
\end{itemize}

\begin{tcolorbox}[breakable,colback=teal!4,colframe=teal!35!black,title={Before you move on}]
What does der Schlüsseldienst mean? (\textquotedblleft{}The locksmith (service)\textquotedblright{}.) Say it once more.
\end{tcolorbox}
\section[aufschließen]{aufschließen — to unlock}

\label{lesson:GE-C445-aufschliessen}



\begin{quote}
Before the new one: say the German for to polish, then the German for the locksmith.
\end{quote}

\subsection*{You'll want to know: aufschließen}
\textbf{aufschließen} — \textquotedblleft{}to unlock\textquotedblright{}.

\begin{quote}
Der Mann schließt die Tür in einer Minute auf.
\end{quote}

\begin{quote}
\emph{The man unlocks the door in one minute.}
\end{quote}

Separable: \emph{ich schließe auf}, \emph{ich habe aufgeschlossen}.

\begin{cousinweb}
\emph{aufschließen} is built from \emph{auf} 'open' and \emph{schließen} 'to lock, to close', inherited from Old High German \emph{sliozan} (EtymDB).
\end{cousinweb}

\begin{grammarlens}[title={What you've built}]
One more word for telling, explaining and describing at B1.
\end{grammarlens}

\subsection*{Your turn}
\begin{itemize}
\raggedright
  \item \emph{Say it:} \emph{aufschließen}
  \item \emph{Say it:} \emph{aufschließen}, once more
  \item \emph{Say it:} a sentence of your own with \emph{aufschließen}
  \item \emph{Recall:} say the German for to polish, then the German for the locksmith, then say \emph{aufschließen} again
\end{itemize}

\begin{tcolorbox}[breakable,colback=teal!4,colframe=teal!35!black,title={Before you move on}]
What does aufschließen mean? (\textquotedblleft{}To unlock\textquotedblright{}.) Say it once more.
\end{tcolorbox}
\section[(dialogue)]{Words from chapters 438–445 — a review}

\label{lesson:GE-R445-review-438-445}



\begin{quote}
Say the German for the old town and the cathedral.
\end{quote}

\subsection*{Your turn}
Say these aloud:

\begin{itemize}
\raggedright
  \item the old town, the cathedral, the bell tower, the town wall, the town gate
  \item the district, the neighbourhood, the suburb, the police station, the barracks
  \item the pavement, the park bench, the street lamp, the cobbles, the fountain
  \item the façade, the shop window, the gutter, the covered market, the pillar
  \item the black ice, to clear, to grit, the snowplough, to slip
  \item the delicatessen, the cake shop, the greengrocer, the chemist, the health food shop
  \item the DIY store, the stationer, the bookshop, the flower shop, to stick
  \item the cobbler, to resole, to polish, the locksmith, to unlock
\end{itemize}

\begin{tcolorbox}[breakable,colback=teal!4,colframe=teal!35!black,title={Before you move on}]
The old town? (\textbf{die Altstadt}.) The cathedral? (\textbf{die Kathedrale}.) The bell tower? (\textbf{der Glockenturm}.) The town wall? (\textbf{die Stadtmauer}.) The town gate? (\textbf{das Stadttor}.) The district? (\textbf{der Stadtteil}.) The neighbourhood? (\textbf{das Viertel}.) The suburb? (\textbf{der Vorort}.) The police station? (\textbf{das Polizeirevier}.) The barracks? (\textbf{die Kaserne}.) The pavement? (\textbf{der Bürgersteig}.) The park bench? (\textbf{die Parkbank}.) The street lamp? (\textbf{die Laterne}.) The cobbles? (\textbf{das Kopfsteinpflaster}.) The fountain? (\textbf{der Brunnen}.) The façade? (\textbf{die Fassade}.) The shop window? (\textbf{das Schaufenster}.) The gutter? (\textbf{die Dachrinne}.) The covered market? (\textbf{die Markthalle}.) The pillar? (\textbf{die Säule}.) The black ice? (\textbf{das Glatteis}.) To clear? (\textbf{räumen}.) To grit? (\textbf{streuen}.) The snowplough? (\textbf{der Schneepflug}.) To slip? (\textbf{rutschen}.) The delicatessen? (\textbf{das Feinkostgeschäft}.) The cake shop? (\textbf{die Konditorei}.) The greengrocer? (\textbf{der Gemüsehändler}.) The chemist? (\textbf{die Drogerie}.) The health food shop? (\textbf{das Reformhaus}.) The DIY store? (\textbf{der Baumarkt}.) The stationer? (\textbf{das Schreibwarengeschäft}.) The bookshop? (\textbf{die Buchhandlung}.) The flower shop? (\textbf{der Blumenladen}.) To stick? (\textbf{kleben}.) The cobbler? (\textbf{der Schuster}.) To resole? (\textbf{besohlen}.) To polish? (\textbf{polieren}.) The locksmith? (\textbf{der Schlüsseldienst}.) To unlock? (\textbf{aufschließen}.)
\end{tcolorbox}
